%%%%%%%%%%%%%%%%%%%%%%%%%%%%%%%%%%%%%%%%%%%%%%%%%%%%%%%%%%%%%%%%%
\chapter*{Introduction}
\addcontentsline{toc}{chapter}{Introduction}
\label{chap:intro}
%%%%%%%%%%%%%%%%%%%%%%%%%%%%%%%%%%%%%%%%%%%%%%%%%%%%%%%%%%%%%%%%%

The rigor of scientific research and the reproducibility of research results are foundational to the validity of scientific findings and the trustworthiness of science as an institution. Yet research rigor and reproducibility remain persistent challenges, particularly for work involving complex data types from heterogeneous sources and long data manipulation pipelines. These challenges are especially acute as data science and artificial intelligence (AI) methods emerge rapidly, and researchers work to seize the opportunities that new methods bring while managing the risks they introduce.

The Data and AI Intensive Research with Rigor and Reproducibility (\DAIR{}) program is a multi-university initiative funded by the National Institute of General Medical Sciences (Award 5R25GM151182). The program equips faculty and technical staff in the biomedical sciences with the skills needed to improve the rigor and reproducibility of their research, and supports them in transferring those skills to their own trainees.

This curriculum document accompanies the \DAIR{} summer bootcamp series. It is organized into seven thematic units covering the principal competencies of the program.

%================================================================
\section*{Curriculum Overview}
\label{sec:intro-overview}
%================================================================

The seven units of this curriculum address the following themes:

\begin{enumerate}
    \item \textbf{Ethics in Biomedical Data Science.} Foundational ethical frameworks, responsible conduct of research, and the application of ethical reasoning to AI and data-driven decision-making in health contexts.
    \item \textbf{Data Management, Representation, and Sharing.} Principles of data lifecycle management, metadata standards, data provenance, and practices that support open and reproducible science.
    \item \textbf{Statistical Foundations.} Core concepts in probability and statistics relevant to biomedical research, with emphasis on rigorous analytical design and appropriate inference.
    \item \textbf{Predictive Modeling.} Design, training, evaluation, and reporting of AI and machine learning models, including considerations of model transparency and fairness.
    \item \textbf{Reproducible Workflows.} Version control, computational environment management, and workflow automation as tools for ensuring that analyses are transparent and reproducible.
    \item \textbf{Meta-Analysis and Assessment Across Studies.} Methods for synthesizing evidence across heterogeneous studies, including systematic reviews and quantitative meta-analysis.
    \item \textbf{Large Language Models.} Principles of generative AI, prompt engineering, and the responsible integration of large language models into biomedical research workflows.
\end{enumerate}

%================================================================
\section*{The 2026 Data Challenge}
\label{sec:intro-challenge}
%================================================================

Each cohort of \DAIR{} participants engages with a data challenge that provides hands-on experience applying the skills developed across the curriculum to a real-world biomedical problem. The 2026 data challenge focuses on \textbf{vital statistics}: projecting the prevalence of underweight newborns and infant mortality by county in Texas through the year 2030. This challenge draws on publicly available natality and mortality data and requires participants to integrate skills from data management, statistical modeling, and reproducible workflow development.

Detailed specifications for the 2026 data challenge, including data sources, analytical requirements, and deliverable formats, are provided in the companion Data Challenge document.

%================================================================
\section*{How to Use This Document}
\label{sec:intro-howto}
%================================================================

Each unit in this curriculum follows a consistent structure. An introductory overview describes the scope and motivation of the unit. Learning objectives state the competencies that participants are expected to develop. The instructional content is organized into sections corresponding to individual sessions or topics. Each unit concludes with assessment instruments that participants complete to demonstrate mastery.

This document is produced in two versions: a classroom version intended for use during instruction, and a datathon version distributed to participants as a reference during the data challenge. Certain sections and annotations differ between versions.

\bigskip
\noindent Participants are encouraged to engage with the material critically, bring questions from their own research contexts, and share knowledge with their peers. The program is designed as a collaborative learning experience across institutions and disciplines.